%%%%%%%%%%  scan idx 307  |  folio 294  %%%%%%%%%%
\SGAPrintPage{294}%
\sgaStableTarget{sga6.e04.para.3-3}3.3. \uline{Relative Grothendieck groups}

\sgaStableTarget{sga6.e04.para.3-3-1}3.3.1. Let \(f:X\longrightarrow Y\) be a locally finite type morphism of schemes. We shall put:
\[
\begin{gathered}
 K^{\bullet}(f)=k(D(f)_{\parf}),\\
 K_{\bullet}(f)=k(D^b(f)_{\coh}),
\end{gathered}
\]
where \(D^b(f)_{\coh}\) denotes the category of complexes pseudo-coherent rel. to \(f\) and with bounded cohomology (III 1.2), and \(D(f)_{\parf}\) the category of complexes of finite perfect amplitude rel. to \(f\) (III 4.2). When no ambiguity is possible, we shall sometimes write \(K^{\bullet}(X/Y)\) (resp. \(K_{\bullet}(X/Y)\)) instead of \(K^{\bullet}(f)\) (resp. \(K_{\bullet}(f)\)).

The groups \(K^{\bullet}(f)\) and \(K_{\bullet}(f)\) do not possess the rich structure of their ``absolute'' analogues \(K^{\bullet}(X)\) and \(K_{\bullet}(X)\), since the tensor product does not in general define a homomorphism \(K^{\bullet}(f)\otimes K^{\bullet}(f)\longrightarrow K^{\bullet}(f)\) (resp. \(K^{\bullet}(f)\otimes K_{\bullet}(f)\longrightarrow K_{\bullet}(f)\)). On the other hand, they satisfy interesting variance properties, which we shall review.

\sgaStableTarget{sga6.e04.para.3-3-2}3.3.2. Let
\[
\begin{tikzcd}
X\arrow[rr,"f"]\arrow[dr,"h"'] &&Y\arrow[dl,"g"]\\
&S&
\end{tikzcd}
\]
be a commutative triangle of locally finite type morphisms of schemes.
\begin{enumerate}[label=3.3.2.\arabic*., leftmargin=4em]
\item If \(g\) is pseudo-coherent, we have \(K_{\bullet}(f)=K_{\bullet}(h)\) (III 1.12). If \(g\) is smooth and \(X\) quasi-compact, we have \(K^{\bullet}(f)=K^{\bullet}(h)\) (III 4.4.1).

\item By (III 1.11 and 4.5), the functor \(\otimes_Y^{\mathrm L}\) defines pairings
\[
 D(f)_{\parf}\times D^b(g)_{\coh}\longrightarrow D^b(h)_{\coh},
\]
\[
 D(f)_{\parf}\times D(g)_{\parf}\longrightarrow D(h)_{\parf},
\]
which induce homomorphisms
\[
 K^{\bullet}(f)\otimes K_{\bullet}(g)\longrightarrow K_{\bullet}(h),
\]
\[
 K^{\bullet}(f)\otimes K^{\bullet}(g)\longrightarrow K^{\bullet}(h).
\]

%%%%%%%%%%  scan idx 308  |  folio 295  %%%%%%%%%%
\SGAPrintPage{295}%
We deduce in particular (by putting \(X=Y\), \(f=\Id\)) that \(K_{\bullet}(g)\) is a \(K^{\bullet}(Y)\)-module. On the other hand, if \(f\) is perfect, multiplication by \(\cl(\underline O_X)\in K^{\bullet}(f)\) defines \(K^{\bullet}(Y)\)-linear homomorphisms
\[
 f^*:K_{\bullet}(Y/S)\longrightarrow K_{\bullet}(X/S),
\]
and
\[
 f^*:K^{\bullet}(Y/S)\longrightarrow K^{\bullet}(X/S).
\]

\item Suppose \(S\) quasi-compact, \(X\) and \(Y\) of finite type over \(S\), and \(f\) proper. Then, subject to the validity of the finiteness conjecture (III 2.1), hence in particular when \(Y\) is noetherian or \(f\) is projective, the functor \(Rf_*\) sends \(D^b(h)_{\coh}\) to \(D^b(g)_{\coh}\), and therefore defines a homomorphism
\[
 f_*:K_{\bullet}(X/S)\longrightarrow K_{\bullet}(Y/S).
\]
If moreover \(f\) is perfect, \(Rf_*\) sends \(D(h)_{\parf}\) to \(D(g)_{\parf}\) (III 4.8), and defines a homomorphism
\[
 f_*:K^{\bullet}(X/S)\longrightarrow K^{\bullet}(Y/S).
\]
From the projection formula (III 3.7) we deduce the formula
\[
 f_*(a x)=a f_*(x)
\]
for \(f\) proper, \(a\in K^{\bullet}(S)\), \(x\in K_{\bullet}(X/S)\) (resp. \(f\) proper and perfect, \(a\in K_{\bullet}(S)\), \(x\in K^{\bullet}(X/S)\)).

\item Suppose \(S\) noetherian, \(X\) and \(Y\) of finite type over \(S\), and \(f\) compactifiable and perfect. Then (III 4.9.1) the functor \(f^!\) sends \(D(g)_{\parf}\) to \(D(h)_{\parf}\), and defines a homomorphism
\[
 f^!:K^{\bullet}(Y/S)\longrightarrow K^{\bullet}(X/S).
\]
\end{enumerate}

\sgaStableTarget{sga6.e04.para.3-3-3}3.3.3. Under the hypotheses of (III 4.7), the functor \(\otimes_S^{\mathrm L}\) defines pairings
\[
 D(f_1)_{\parf}\times D^b(f_2)_{\coh}\longrightarrow D^b(f)_{\coh},
\]
\[
 D(f_1)_{\parf}\times D(f_2)_{\parf}\longrightarrow D(f)_{\parf},
\]
which induce homomorphisms
\[
 K^{\bullet}(f_1)\otimes K_{\bullet}(f_2)\longrightarrow K_{\bullet}(f),
\]
\[
 K^{\bullet}(f_1)\otimes K^{\bullet}(f_2)\longrightarrow K^{\bullet}(f).
\]

%%%%%%%%%%  scan idx 309  |  folio 296  %%%%%%%%%%
\SGAPrintPage{296}%
(Of course, for \(X_i=Y_i=S\), we recover the \(K^{\bullet}(S)\)-module structure of \(K_{\bullet}(S)\).)

If \(X\) is flat over \(S\), we deduce, for \(n\geq1\), homomorphisms
\[
 K^{\bullet}(X/S)^{\otimes n}\longrightarrow K^{\bullet}\bigl((X/S)^n\bigr),
\]
which should in a sense replace the absent ring structure on \(K^{\bullet}(X/S)\).

\sgaStableTarget{sga6.e04.bibliography}\subsection*{BIBLIOGRAPHY}
\begin{enumerate}[label={[\arabic*]}]
\item Deligne, P., papiers secrets.
\item Grauert, H., Ein Theorem der Analytischen Garbentheorie \dots, Pub. Math. IHES, n\textsuperscript{o}5 (1960).
\item Grothendieck, A., Théorèmes de finitude pour la cohomologie des faisceaux, Bull. Soc. Math. France, 84, 1956, p.1-7.
\item[{[V]}] Verdier, J.-L., Catégories dérivées, IHES (1964), miméographié.
\end{enumerate}


% SGA_NATIVE_CONTINUATION_296
\providecommand{\Chat}{\widehat C}
\providecommand{\Ghat}{\widehat G}
\providecommand{\Gg}{\widehat G_g}
\providecommand{\Ga}{\widehat G_a}
\providecommand{\Grhat}{\widehat{\operatorname{Gr}}}
\providecommand{\Abfil}{\mathrm{Ab\,fil}}
\providecommand{\Kalgfil}{K\text{-}\mathrm{Alg.fil.}}
\providecommand{\Hom}{\operatorname{Hom}}
\providecommand{\Sym}{\mathfrak S}
\providecommand{\id}{\operatorname{Id}}
\providecommand{\Log}{\operatorname{Log}}
\providecommand{\Toddc}{\operatorname{Todd}}
\providecommand{\Ob}{\operatorname{Ob}}
\providecommand{\cG}{\mathcal G}
\providecommand{\cCzero}{\mathcal C_0}
\providecommand{\Fil}{\operatorname{Fil}}



% SGA_REVIEWED_FRENCH_CONTINUATION_296
% Expose V macros
\providecommand{\Chat}{\widehat C}
\providecommand{\Ghat}{\widehat G}
\providecommand{\Gg}{\widehat G_g}
\providecommand{\Ga}{\widehat G_a}
\providecommand{\Grhat}{\widehat{\operatorname{Gr}}}
\providecommand{\Abfil}{\underline{\mathrm{Ab}}\,\underline{\mathrm{fil}}}
\providecommand{\Kalgfil}{K\text{-}\mathrm{Alg.fil.}}
\providecommand{\Hom}{\operatorname{Hom}}
\providecommand{\Sym}{\mathfrak S}
\renewcommand{\id}{\operatorname{Id}}
\providecommand{\Log}{\operatorname{Log}}
\providecommand{\Toddc}{\operatorname{Todd}}
\renewcommand{\Ob}{\operatorname{Ob}}
\providecommand{\cG}{\mathcal G}
\providecommand{\cCzero}{\mathcal C_0}
\providecommand{\Fil}{\operatorname{Fil}}



% SGA_SOURCE_MACRO_RESTORATION_296: native provide must not shadow the two reviewed word rules.
\renewcommand{\Abfil}{\underline{\mathrm{Ab}}\,\underline{\mathrm{fil}}}

